\begin{statusdone}
레이더--카메라 거리 융합이 실시간으로 동작한다. 본 시스템의 핵심 성과다.
\end{statusdone}

이 장은 본 프로젝트의 핵심이다. 레이더의 정확한 거리와 카메라의 의미 정보를
방위각으로 대응시켜, ``각 객체까지의 실제 거리''를 추정한다.

\chapter{원리 및 설계 근거}

\section{레이더 검출점의 화면 투영}
레이더 검출 $(r, \theta_{rad})$를 카메라 화면의 $x$ 좌표로 투영한다. 먼저 레이더
좌표계에서 직교 좌표로,
\[
  x_r = r\sin\theta_{rad},\qquad z_r = r\cos\theta_{rad},
\]
외부 캘리브(yaw 회전 + 평행이동 $t_x, t_z$)로 카메라 좌표계로 변환,
\[
  x_c = \cos\psi\, x_r + \sin\psi\, z_r + t_x,\qquad
  z_c = -\sin\psi\, x_r + \cos\psi\, z_r + t_z,
\]
핀홀 투영으로 화면 좌표,
\[
  s_x = c_x + f_x \frac{x_c}{z_c}.
\]
($\psi$: yaw 오프셋, $f_x, c_x$: 카메라 내부 파라미터.)

\section{외부 캘리브레이션 — 방위각 대응}
yaw 오프셋 $\psi$를 추정하기 위해, 카메라가 본 방위각과 레이더 방위각을 대응시킨다.
객체 bbox 중심 $b_{cx}$에 대해
\[
  \theta_{cam} = \arctan\!\frac{b_{cx}-c_x}{f_x},
\]
여러 표본에서 $\psi = \mathrm{median}(\theta_{cam}-\theta_{rad})$로 구한다.

\begin{designbox}{2단계 캘리브 (near $\to$ fine)}
가까운 거리(1$\sim$2\,m)에서 사람이 좌우로 움직이며 표본을 모아 yaw를 먼저 수렴
시키고(near), 수렴하면 먼 거리(최대 6\,m)를 추가해 정밀화한다(fine). 중복 표본
(거리/방위각이 비슷한 점)은 거른다. 수렴 시 \texttt{config/calib\_radar\_cam.json}에
자동 저장되어 융합에서 자동 로드된다. 실측 yaw $\approx 0.8^\circ$, 잔차 RMS
$\approx 5.5^\circ$.
\end{designbox}

\section{방위각 프로파일 매칭}
객체의 화면 가로 범위(마스크가 있으면 마스크의 열 투영, 없으면 bbox)를
프로파일로 만들고, 투영된 레이더 점 $s_x$가 그 프로파일 안에 들어오는지로 1차
매칭한다(\texttt{\_in\_profile}). 마스크 열 투영을 쓰면 직사각형 bbox보다 객체
실제 윤곽에 가깝게 점을 가른다.

\section{거리 게이트 — expected 기반}
한 객체 프로파일에 여러 점이 들어올 수 있으므로, 깊이 추정으로 만든
\emph{기대 거리}로 2차 게이트를 건다.
\[
  \text{expected} = \text{da3\_0}\times\text{ratio},\quad
  \text{gate} = \max(1.2,\ \text{expected}\times 0.5),
\]
$|r-\text{expected}|\le\text{gate}$인 점만 채택하고, 여러 객체에 걸리면 기대 거리에
가장 가까운 객체로 배정한다(\texttt{match\_all}). 채택된 점들의 range/azimuth/velocity
중앙값, SNR/power 최댓값을 객체 측정값으로 집계한다.

\section{신뢰 기반 거리 결정}
\begin{itemize}
  \item \textbf{SNR 게이트}: \texttt{SNR\_MIN=4.0}. 레이더 검출이 충분히 강하면
        레이더 거리를 채택, 약하면(또는 미검출) 깊이 추정으로 대체.
  \item \textbf{깊이 추정 대체(est)}: $\text{거리} = \text{da3\_0}\times\text{scale}$.
\end{itemize}

\begin{designbox}{스케일 신뢰 게이트 + 이상치 거부}
객체별 레이더/깊이 비율(ratio)을 EMA($0.7\cdot\text{old}+0.3\cdot\text{new}$)로 학습하되,
새 비율이 기존/전역 대비 $\times2$ 또는 $\div2$ 범위를 벗어나면 \emph{이상치로 거부}한다.
또한 누적 학습 횟수 \texttt{ratio\_n}을 세어, $\ge 5$회 일관되게 잡힌 객체만 자기 ratio를
신뢰하고, 한 번도/짧게만 잡힌 객체는 \emph{신뢰 객체들의 전역 중앙값}을 적용한다.
이로써 레이더가 튀어도 스케일이 오염되지 않고, 부정확한 객체가 잘못된 자기 스케일을
쓰는 것을 막는다.
\end{designbox}

\begin{designbox}{속도(v) 적응 거리 EMA}
정지 물체는 도플러 $v\approx 0$이 물리적으로 robust한데, 거리 $r$만 약검출
노이즈로 널뛰는 경우가 많다(관측: 정지 의자가 $v=+0.00$인데 $r$이 $3.2\sim4.8$\,m
변동). 이를 이용해 객체별 거리 상태 $r_{sm}$을 다음으로 평활한다.
\[
  \alpha = \begin{cases}0.50 & |v|\ge 0.1 \ (\text{동적, 추종})\\
                         0.15 & |v|<0.1 \ (\text{정지, 노이즈 흡수})\end{cases},\qquad
  r_{sm}\leftarrow r_{sm}+\alpha\,(r_{meas}-r_{sm}).
\]
정지 물체는 강하게 평활해 안정시키고, 움직이는 물체는 빠르게 따라간다.
\end{designbox}

\chapter{구현 및 디렉토리 구조}

\section{\texttt{src/utils/radar\_cam\_calib.py}}
\texttt{calibrate\_loop} — \texttt{live\_overlay.json}의 객체(range, az, bbox, n)를
폴링하며 표본을 모아 yaw를 추정. near/fine 2단계, 수렴 시 JSON 저장.

\section{\texttt{src/objects/radar\_fusion.py}}
\begin{itemize}
  \item \texttt{radar\_to\_x} — 위 투영식(외부 캘리브 자동 로드).
  \item \texttt{\_azimuth\_profile} / \texttt{\_in\_profile} — 마스크 열 투영 프로파일.
  \item \texttt{match\_all} — expected 게이트로 점→객체 배정.
  \item \texttt{\_aggregate} — range/az/v 중앙값 + SNR/power 최댓값 집계.
  \item \texttt{load\_latest\_targets} — 카메라 프레임 시각과 가장 가까운 레이더
        프레임 선택(시간 허용오차), 비어 있으면 직전 비어있지 않은 값 유지.
\end{itemize}

\section{\texttt{main\_r.py} 융합 부분}
\texttt{da3\_0}(첫 깊이 고정) → \texttt{expected} 계산 → \texttt{match\_all} →
SNR/ratio 학습(\texttt{ratio\_n}, 이상치 거부) → \texttt{\_obj\_scale}(신뢰 게이트) →
실측이면 v 적응 EMA, 약검출이면 est. 결과를 \texttt{live\_overlay.json}에 기록.

\begin{figure}[H]
  \centering
  % figures/viewer_2d.png 에 2D 뷰어 스크린샷을 넣으세요
  \includegraphics[width=0.9\textwidth]{viewer_2d}
  \caption{2D 뷰어 — 각 객체에 융합 거리가 라벨로 표시된다. 레이더 실측(R),
  단안 깊이(D), 최종 융합 거리를 \texttt{dist} 모드로 비교할 수 있다.}
  \label{fig:viewer2d}
\end{figure}

\chapter{문제 해결 과정}

\section{정지 객체 거리 고착 / 널뛰기}
\textbf{증상}: 정지 의자의 거리가 깊이 추정값에 고착되거나, 반대로 $3\sim5$\,m로
널뛰었다.\\
\textbf{분석}: 콘솔에 $v$를 함께 출력해 확인 — 정지 물체는 $v\approx0$로 안정인데
range만 노이즈였다(레이더 range/v는 같은 검출점에서 나오지만, 정지 물체의
도플러 0은 robust).\\
\textbf{해결}: $|v|$ 적응 EMA(위 설계 박스). 정지=강한 평활, 동적=추종.

\section{스케일 오염}
\textbf{증상}: 한 번 튄 레이더 값이 객체 스케일을 망가뜨림.\\
\textbf{해결}: 이상치 거부($\times2/\div2$) + 신뢰 게이트(\texttt{ratio\_n}$\ge5$) +
신뢰 객체만으로 전역 중앙값 산출.

\section{거리 게이트의 부작용 정리}
초기엔 고정 거리 게이트가 정지 객체까지 흔들었다. 깊이 추정 기반
\emph{expected} 게이트로 바꿔, 물체마다 적응적인 허용 범위를 갖게 했다.
